\documentclass[a4paper,11pt,fleqn]{article}%fleqn pour aligner les equations à gauche
\usepackage{../../Styles/Cours}


\newcommand{\chnum}{Chapitre 16}
\newcommand{\chtitle}{Caractéristiques d'une lentille mince convergente}
\newcommand{\dtype}{TP}

\begin{document}

\maketitle
\noindent \fbox{\begin{minipage}{.99\linewidth}
    Afin d’étudier les propriétés d’une lentille inconnue on dispose d’un dispositif, nommé le banc optique, permettant de mesurer les propriétés caractéristiques d’une lentille.
Un objet lumineux AB (ici une lettre) est placé sur une graduation du banc. La lentille est placée à une position déterminée et on place l’écran de manière à obtenir une image la plus nette possible.
Le grandissement $\gamma$ (gamma) est le rapport entre la taille de l’image notée A’B’ et la taille de l’objet notée AB. On a donc : $$\gamma=\dfrac{A'B'}{AB}$$ \begin{center}Si $\gamma>1$ l’image est agrandie, si $\gamma<1$l’image est rétrécie.\end{center}
\end{minipage}}

\textbf{Procéder à des mesures :}
\begin{enumerate}
    \item Faire un schéma du banc optique (On choisira de placer le $0$ de l’axe horizontal à la position de la lentille)
    \item Prendre la mesure d’une lettre servant d’objet $AB$.
    \item Placer la lentilles aux différentes distances $OA$ du tableau puis déplacer l’écran pour avoir une image nette. Compléter le tableau.
\end{enumerate}
\begin{center}
    \begin{tabular}{|>{\centering\columncolor{gray}}m{.2\textwidth}|>{\centering}m{.1\textwidth}|>{\centering}m{.1\textwidth}|>{\centering}m{.1\textwidth}|>{\centering\arraybackslash}m{.1\textwidth}|}
    \hline Distance Objet - Lentille (m) & 0,60 & 0,50 & 0,40 & 0,30 \\
    \hline Distance Lentille - écran (m)& & & & \\
    \hline Hauteur de l'image $AB$ (m)& & & & \\
    \hline Image droite ou renversée ?& & & & \\
    \hline Image agrandie ou rétrécie ?& & & & \\
    \hline
\end{tabular}
\end{center}
\begin{enumerate}[resume]
    \item Comment évoluent la distance $OA’$ et la taille de l’image $A’B’$ lorsque l’on augmente la distance $OA$
    \item Que se passe t’il si la distance entre l’objet et la lentille est inférieure à la distance focale (ici $f = \SI{25}{cm}$) ? (faire l’expérience).
    \item A partir des différentes mesures, calculer pour chaque distance le grandissement et le rapport entre $OA$ et $OA’$.
\end{enumerate}
\begin{center}
    \begin{tabular}{|>{\centering\columncolor{gray}}m{.2\textwidth}|>{\centering}m{.1\textwidth}|>{\centering}m{.1\textwidth}|>{\centering}m{.1\textwidth}|>{\centering\arraybackslash}m{.1\textwidth}|}
    %\hline Distance Objet - Lentille (m) & 0,60 & 0,50 & 0,45 & 0,40 \\
    \hline $\gamma =\dfrac{A'B'}{AB}$& & & & \\
    \hline $\dfrac{OA'}{OA}$& & & & \\
    %\hline Image droite ou renversée ?& & & & \\
    %\hline Image agrandie ou rétrécie ?& & & & \\
    \hline
\end{tabular}
\end{center}
\begin{enumerate}[resume]
    \item Comparer les valeurs obtenues pour le grandissement et pour $\dfrac{OA'}{OA}$
    \item En déduire une autre manière de calculer le grandissement.
    \item Représenter sur deux schémas différents les situations observée pour $OA = \SI{40}{cm}$ et pour $OA = \SI{20}{cm}$, en utilisant le modèle des rayons lumineux. Prendre pour échelle \SI{5}{cm} réels = \SI{1}{cm} sur le schéma.
\end{enumerate}

\end{document}
